% Capítulo de Metodología

El proyecto evolucionó a través de tres prototipos incrementales: P0 (web scraping directo, fallido), P1 (RSS y clustering básico) y P2 (triple fuente con detector especializado). Cada prototipo representa un ciclo de desarrollo donde los fallos guiaron el diseño del siguiente.

\section{Fase 0: Web Scraping Directo (Fallido)}
\label{sec:prototipo0}

\subsection{Objetivo}
Validar la extracción directa de noticias mediante scraping tradicional con requests y BeautifulSoup4.

\subsection{Implementación}
Aplicación Flask que codificaba URLs de medios, descargaba HTML, parseaba contenido y buscaba "feminicidio".

\subsection{Fallo Completo}
Tres razones bloquearon el prototipo:

\begin{description}
    \item[Bloqueo Anti-Bot:] HTTP 403/406 por ausencia de headers válidos de navegador.
    \item[Contenido Dinámico:] JavaScript rendering que \texttt{requests} no ejecuta.
    \item[Cloudflare:] Desafíos JavaScript imposibles de superar con herramientas básicas.
\end{description}

\textbf{Resultado:} 0 noticias de 10 fuentes intentadas.

\subsection{Lecciones}
Abandono del scraping HTML directo. Migración a fuentes estructuradas (RSS, APIs) como estrategia principal.

\section{Fase 1: Sistema Básico Funcional}
\label{sec:prototipo1}

\subsection{Objetivo}
Sistema funcional mínimo (sept 24 - oct 2, 2025) que resuelva bloqueos del P0.

\subsection{Solución Implementada: Feeds RSS}
La estrategia clave fue abandonar el web scraping directo y adoptar \textbf{Feeds RSS/Atom}, que son:
\begin{itemize}
    \item \textbf{Estructurados:} Formato XML estandarizado, fácil de parsear con \texttt{feedparser}.
    \item \textbf{Públicos:} Diseñados para ser consumidos por terceros, sin bloqueos anti-bot.
    \item \textbf{Actualizados:} Los medios publican automáticamente sus últimas noticias.
\end{itemize}

\subsection{Implementación del Prototipo 1}

\subsubsection{Arquitectura}
Sistema monolítico Flask con tres componentes:
\begin{enumerate}
    \item \texttt{data\_collector.py}: Recolector RSS que consulta 8 fuentes mexicanas cada 6 horas.
    \item \texttt{simplified\_analyzer.py}: Flujo ML básico con 4 etapas (TF-IDF → K-Means → LDA → Similarity).
    \item \texttt{app.py}: API REST Flask con 3 endpoints (\texttt{/api/noticias}, \texttt{/api/search}, \texttt{/api/analyze}).
\end{enumerate}

\subsubsection{Fuentes RSS Configuradas}
\begin{table}[h]
\centering
\caption{Fuentes RSS del Prototipo 1}
\begin{tabular}{|l|l|}
\hline
\textbf{Medio} & \textbf{Sección} \\ \hline
La Jornada & Política \\
Proceso & General \\
Aristegui Noticias & Nacional \\
Animal Político & Noticias \\
Sin Embargo & MX \\
Forbes México & Actualidad \\
El Sol de México & México \\
El Financiero & Nacional \\ \hline
\end{tabular}
\end{table}

\subsubsection{Flujo ML Básico}
\begin{description}
    \item[Paso 1 - Vectorización:] \texttt{TfidfVectorizer} con \texttt{max\_features=1000} para convertir texto a matriz numérica (reducido vs. 3000 actual).
    \item[Paso 2 - Clustering:] \texttt{KMeans} con \texttt{n\_clusters=4} para agrupar noticias del mismo evento (fijo vs. adaptativo DBSCAN).
    \item[Paso 3 - Tópicos:] \texttt{LatentDirichletAllocation} con \texttt{n\_components=5} para extraer temas generales.
    \item[Paso 4 - Similitud:] Cosine similarity con threshold $\geq 0.80$ para detectar duplicados (más estricto vs. 0.75 actual).
\end{description}

\subsubsection{Detección NNA Rudimentaria}
Se implementó una función \texttt{detect\_children\_mentions()} basada en regex simple que buscaba palabras clave como: 'hijo', 'hija', 'niña', 'niño', 'menor', 'huérfano', 'bebé'. Sin ponderación de contexto ni exclusión de falsos positivos.

\subsection{Resultados del Prototipo 1}

\subsubsection{Métricas Cuantitativas}
\begin{itemize}
    \item \textbf{Noticias recolectadas:} 96 noticias en ejecución típica (promedio de 12 por fuente RSS).
    \item \textbf{Tasa de detección NNA:} 28.1\% (27/96 noticias), pero con alta tasa de falsos positivos ($\sim$40\%).
    \item \textbf{Clusters generados:} 4 clusters fijos por K-Means (independiente de la distribución real de datos).
    \item \textbf{Tópicos identificados:} 5 tópicos, pero con interpretabilidad baja debido a vocabulario limitado (1000 features).
    \item \textbf{Tiempo de ejecución:} $\sim$45 segundos por análisis completo.
\end{itemize}

\subsection{Limitaciones Críticas del Prototipo 1}

\begin{enumerate}
    \item \textbf{Baja precisión semántica:} El detector regex marcaba como positivas noticias sobre "leyes de protección infantil", "campañas de prevención" o "estadísticas de violencia", generando falsos positivos del 40\%.
    
    \item \textbf{Cobertura limitada:} La estrategia de "solo RSS" cubría únicamente 8 fuentes, dejando fuera medios regionales importantes y fuentes especializadas en temas de género (como CIMAC Noticias). Volumen: $\sim$70-100 noticias/ejecución.
    
    \item \textbf{Clustering rígido:} K-Means con $k=4$ fijo no se adapta a la realidad de los datos. Los feminicidios son eventos únicos dispersos, pero el algoritmo fuerza 4 agrupaciones aunque no exista similitud real. Silhouette Score: 0.23 (clustering pobre).
    
    \item \textbf{Almacenamiento no escalable:} CSV como única persistencia genera:
    \begin{itemize}
        \item Problemas de encoding con caracteres especiales (UTF-8 vs. UTF-8-BOM en Windows).
        \item Imposibilidad de consultas complejas (búsqueda por rango de fechas, filtros múltiples).
        \item Pérdida de datos al sobrescribir archivos en ejecuciones concurrentes.
    \end{itemize}
    
    \item \textbf{Sin containerización:} Ejecución manual requiere configurar entorno Python en cada máquina, generando inconsistencias (versiones de librerías, paths relativos vs. absolutos).
\end{enumerate}

\subsection{Conclusiones del Prototipo 1}
\textbf{Éxito:} Validamos que RSS resuelve el problema de bloqueo y que el Flujo ML básico (TF-IDF → Clustering → Tópicos) es técnicamente viable.

\textbf{Fracaso:} La precisión del sistema es insuficiente para uso real (28.1\% detección con 40\% falsos positivos = solo 17\% casos verdaderos). Las limitaciones del P1 definen los requisitos del P2.

\section{Fase 2: Prototipo 2 - Sistema Robusto de Producción}
\label{sec:prototipo2}

A partir de las lecciones críticas del P1, diseñamos el Prototipo 2 (noviembre 11-19, 2025, rama \texttt{pruebas3}) para resolver los cinco problemas identificados y construir un sistema listo para validación en entorno real. Este prototipo representa la culminación del TT1 y establece la base arquitectónica para el TT2.

\subsection{Requerimientos Implementados}
\label{subsec:p2_requerimientos}

Este prototipo implementa la totalidad de los requerimientos funcionales y no funcionales críticos del TT1:

\begin{description}
    \item[RF-01, RF-02 (Recolección Robusta):] Resueltos al evolucionar de 8 RSS a una arquitectura de \textbf{triple fuente}: RSS (10 feeds) + Google News API (5 queries) + Historical Scraper (6 meses retrospectivos). Implementado en \texttt{data\_collector.py}, \texttt{feminicide\_detector.py} y \texttt{historical\_scraper.py}.
    
    \item[RF-04, RF-10 (Análisis y Clasificación Avanzados):] Resueltos con el Flujo ML optimizado en \texttt{simplified\_analyzer.py}: TF-IDF (3000 features) → LDA (8 topics) → DBSCAN (clustering adaptativo) → Similarity (75\% threshold). Clasificación de prioridad: ALTA ($\geq$70\%), MEDIA (40-69\%), BAJA (20-39\%).
    
    \item[RF-03 (Detector Especializado):] Implementado \texttt{FeminicideDetector} con 40+ patrones de feminicidio, 15+ patrones NNA y 17 patrones de exclusión (legislación, estadísticas, campañas). Tasa de detección: 44.8\% (126/281 noticias).
    
    \item[RF-06, RF-16 (Búsqueda Semántica):] Implementada con diccionario de sinónimos especializado (\texttt{synonym\_dictionary.py}) con 127 términos agrupados en 18 categorías semánticas. API endpoint: \texttt{/api/search?q=<query>}.
    
    \item[RNF-02, RNF-05 (Escalabilidad, Mantenibilidad):] Resueltos con arquitectura de microservicios Docker: contenedor \texttt{nna-analyzer} (worker) + contenedor \texttt{nna-webapp} (Flask API + Dashboard). Orquestación: \texttt{docker-compose.yml}.
    
    \item[RNF-03 (Persistencia):] Implementada estrategia multi-archivo: \texttt{noticias\_raw.csv} (datos crudos), \texttt{noticias\_analyzed\_simplified.csv} (datos procesados), \texttt{reporte\_ml.json} (métricas), \texttt{clusters\_info.csv} (clustering). Encoding: UTF-8-sig para compatibilidad Windows/Excel.
\end{description}

\subsection{Stack Tecnológico Completo}
\label{subsec:p2_stack}

\begin{description}
    \item[Python 3.11] Lenguaje principal del sistema.
    \item[Pandas 2.0+] Manipulación de datasets (DataFrames con 280+ filas × 15 columnas).
    \item[Scikit-learn 1.3+] Flujo ML: \texttt{TfidfVectorizer}, \texttt{LatentDirichletAllocation}, \texttt{DBSCAN}, \texttt{cosine\_similarity}.
    \item[Requests + BeautifulSoup4] Recolección RSS y parsing XML/HTML.
    \item[Google News API (gnews)] Consultas programáticas con rate limiting (2-5s delay, checkpoints cada 30 requests).
    \item[Flask 3.0] Framework web para API REST (6 endpoints) y dashboard Bootstrap 5.
    \item[Docker + Docker Compose] Containerización y orquestación de microservicios.
    \item[APScheduler] Ejecución programada del worker (scheduler de 24 horas).
\end{description}

\subsection{Arquitectura del Sistema}
\label{subsec:p2_arquitectura}

La arquitectura del P2 es una evolución radical respecto al P1. Migramos de un monolito Flask a un sistema distribuido de microservicios desacoplados:

\begin{figure}[h]
    \centering
    \framebox(12cm, 7cm){[Diagrama de Arquitectura de Microservicios]}
    \caption{Arquitectura de microservicios del Prototipo 2 con triple fuente de datos.}
    \label{fig:arquitectura_p2}
\end{figure}

\subsubsection{Componentes del Sistema}

\begin{enumerate}
    \item \textbf{nna-analyzer (Backend Worker):}
    \begin{itemize}
        \item Ejecuta el Flujo completo de recolección y análisis ML.
        \item Scheduler: APScheduler con intervalo de 24 horas (configurable).
        \item Output: Genera 5 archivos CSV/JSON en volumen compartido \texttt{./data}.
        \item Sin interfaz web (headless service).
    \end{itemize}
    
    \item \textbf{nna-webapp (Frontend + API):}
    \begin{itemize}
        \item Framework: Flask 3.0 con templates Jinja2 + Bootstrap 5.
        \item API REST: 6 endpoints (\texttt{/api/stats}, \texttt{/api/noticias}, \texttt{/api/search}, \texttt{/api/analyze}, \texttt{/api/export/csv}, \texttt{/api/health}).
        \item Dashboard interactivo: 4 secciones (estadísticas globales, gráficos ML, tabla de noticias paginada, búsqueda avanzada).
        \item Puerto: 5000 (mapeado a host).
    \end{itemize}
    
    \item \textbf{Volumen Compartido Docker (\texttt{./data}):}
    \begin{itemize}
        \item Comunicación asíncrona entre microservicios.
        \item El worker escribe, la webapp lee (patrón producer-consumer).
        \item Persistencia: Los datos sobreviven al reinicio de contenedores.
    \end{itemize}
\end{enumerate}

\subsubsection{Ventajas de la Arquitectura de Microservicios}

\begin{itemize}
    \item \textbf{Escalabilidad horizontal:} Podemos ejecutar múltiples workers en paralelo para aumentar frecuencia de recolección (cada 6 horas en lugar de 24).
    \item \textbf{Despliegue independiente:} Actualizar el dashboard no requiere reiniciar el worker (y viceversa).
    \item \textbf{Tolerancia a fallos:} Si la webapp falla, el worker continúa recolectando datos.
    \item \textbf{Portabilidad:} \texttt{docker-compose up} funciona idénticamente en Windows, Linux y macOS.
\end{itemize}

\subsection{Implementación del Flujo de Recolección y Análisis}
\label{subsec:p2_pipeline}

El P2 implementa un Flujo completo de 7 etapas, representando un incremento del 75\% en complejidad respecto al P1 (4 etapas):

\subsubsection{Etapa 1: Recolección Multi-Fuente}

\textbf{Innovación clave:} Arquitectura de triple fuente que combina volumen (RSS), precisión (Google News) y profundidad histórica (Historical Scraper).

\paragraph{Fuente 1: RSS Feeds (70 noticias/ejecución)}
\begin{itemize}
    \item \textbf{Configuración:} 10 feeds de medios mexicanos (vs. 8 en P1).
    \item \textbf{Implementación:} \texttt{data\_collector.py} con librería \texttt{feedparser}.
    \item \textbf{Ventajas:} Actualizaciones automáticas cada 6 horas, sin bloqueos, formato XML estructurado.
    \item \textbf{Desventaja:} Volumen limitado (promedio 7 noticias/feed).
\end{itemize}

\paragraph{Fuente 2: Google News API (150 noticias/ejecución)}
\begin{itemize}
    \item \textbf{Queries configuradas:} 5 búsquedas especializadas ('feminicidio México', 'asesinato mujer', 'violencia género', 'crimen mujer', 'homicidio doloso mujer').
    \item \textbf{Rate Limiting:} Delay aleatorio 2-5 segundos entre requests, checkpoints cada 30 peticiones para evitar HTTP 429 (Too Many Requests).
    \item \textbf{Ventajas:} Cobertura nacional, incluye medios regionales sin RSS.
    \item \textbf{Limitación:} API gratuita limita a 100 resultados/query/día.
\end{itemize}

\paragraph{Fuente 3: Historical Scraper (250 noticias/ejecución)}
\begin{itemize}
    \item \textbf{Período:} Recolección retrospectiva de 6 meses (mayo-noviembre 2025).
    \item \textbf{Implementación:} \texttt{historical\_scraper.py} con búsquedas paginadas en Google News.
    \item \textbf{Ventajas:} Proporciona contexto histórico para análisis de tendencias temporales.
    \item \textbf{Ejecución:} Una sola vez al inicializar el sistema (no repetida en ejecuciones programadas).
\end{itemize}

\textbf{Volumen total:} $\sim$470 noticias raw $\rightarrow$ $\sim$280-300 únicas después de deduplicación (incremento de 300\% vs. P1).

\subsubsection{Etapa 2: Deduplicación Avanzada}

Algoritmo de similitud de coseno con threshold calibrado en 75\% (vs. 80\% en P1):

\begin{equation}
similarity(d_1, d_2) = \frac{\mathbf{v}_1 \cdot \mathbf{v}_2}{||\mathbf{v}_1|| \cdot ||\mathbf{v}_2||} \geq 0.75
\end{equation}

donde $\mathbf{v}_i$ es el vector TF-IDF del documento $i$.

\textbf{Resultado:} Detecta $\sim$6 duplicados por ejecución (1.3\% de falsos positivos eliminados). Threshold 75\% es más permisivo que 80\% porque diferentes medios usan vocabulario distinto para el mismo evento.

\subsubsection{Etapa 3: Detección Especializada de Feminicidios}

\textbf{Componente nuevo:} Clase \texttt{FeminicideDetector} en \texttt{feminicide\_detector.py}.

\paragraph{Sistema de Patrones Regex (72 patrones totales)}

\begin{description}
    \item[40+ patrones de feminicidio:] 'feminicidio', 'asesinato.*mujer', 'homicidio doloso.*femenino', 'víctima.*violencia género', 'crimen.*machista', etc.
    
    \item[15+ patrones NNA:] 'hijo.*huérfano', 'menor.*edad', 'niña.*viuda', 'bebé.*madre asesinada', 'adolescente.*orfandad', etc.
    
    \item[17 patrones de exclusión:] 'ley.*protección', 'campaña.*prevención', 'estadística', 'capacitación.*jueces', 'conferencia.*prensa', 'trata.*personas' (para eliminar falsos positivos).
\end{description}

\paragraph{Cálculo de Confianza}

\begin{equation}
confidence = \frac{matches_{feminicidio} + matches_{NNA}}{total\_patrones} \times 100
\end{equation}

\begin{equation}
prioridad = 
\begin{cases}
ALTA & \text{si } confidence \geq 70\% \\
MEDIA & \text{si } 40\% \leq confidence < 70\% \\
BAJA & \text{si } 20\% \leq confidence < 40\% \\
IRRELEVANTE & \text{si } confidence < 20\%
\end{cases}
\end{equation}

\textbf{Resultados:} Tasa de detección 44.8\% (126/281 noticias), vs. 28.1\% en P1 (incremento de 59\%). Tasa de falsos positivos reducida a $\sim$10\% (vs. 40\% en P1).

\subsubsection{Etapa 4: Vectorización TF-IDF}

Conversión de texto a matriz numérica con \texttt{TfidfVectorizer}:

\begin{lstlisting}[language=Python, caption={Configuración TF-IDF del Prototipo 2}]
vectorizer = TfidfVectorizer(
    max_features=3000,      # vs. 1000 en P1 (+200% vocabulario)
    ngram_range=(1, 2),     # unigramas y bigramas
    min_df=1,               # minima frecuencia documento
    max_df=0.8,             # maxima (elimina stopwords)
    encoding='utf-8-sig'    # compatibilidad Windows
)
\end{lstlisting}

\textbf{Output:} Matriz de dimensión $126 \times 3000$ (vs. $96 \times 1000$ en P1).

\subsubsection{Etapa 5: Modelado de Tópicos con LDA}

Latent Dirichlet Allocation para identificar temas latentes:

\begin{lstlisting}[language=Python, caption={Configuración LDA del Prototipo 2}]
lda = LatentDirichletAllocation(
    n_components=8,     # vs. 5 en P1 (+60% granularidad)
    random_state=42,
    max_iter=10
)
\end{lstlisting}

\textbf{Interpretación:} Cada noticia se distribuye probabilísticamente sobre 8 tópicos (ej: 40\% Tópico 1 [Violencia Doméstica] + 35\% Tópico 3 [Justicia] + 25\% Tópico 5 [Impunidad]).

\subsubsection{Etapa 6: Clustering Adaptativo con DBSCAN}

\textbf{Cambio crítico:} Reemplazo de K-Means por DBSCAN.

\begin{lstlisting}[language=Python, caption={Configuración DBSCAN del Prototipo 2}]
dbscan = DBSCAN(
    eps=0.8,            # radio de vecindad
    min_samples=2,      # minimo de puntos por cluster
    metric='cosine'     # distancia semantica
)
\end{lstlisting}

\paragraph{Justificación del Cambio}

\begin{table}[h]
\centering
\caption{Comparativa K-Means vs. DBSCAN en Prototipo 2}
\begin{tabular}{|l|c|c|}
\hline
\textbf{Métrica} & \textbf{K-Means (P1)} & \textbf{DBSCAN (P2)} \\ \hline
Silhouette Score & 0.23 & 0.48 \\ \hline
Clusters detectados & 4 (fijo) & 3-6 (adaptativo) \\ \hline
Outliers & 0 (todos forzados) & 81.4\% \\ \hline
Interpretabilidad & Baja & Alta \\ \hline
\end{tabular}
\end{table}

\textbf{} Los feminicidios son eventos únicos dispersos, no agrupaciones densas. El 81.4\% de outliers es \textbf{correcto} porque representa casos sin duplicados mediáticos, mientras que los 3-6 clusters identifican eventos con cobertura masiva (ej: feminicidio de Debanhi Escobar con 15+ noticias).

\subsubsection{Etapa 7: Búsqueda Semántica con Sinónimos}

Expansión de consultas usando tesauro especializado (\texttt{synonym\_dictionary.py}):

\begin{itemize}
    \item \textbf{127 términos} agrupados en \textbf{18 categorías semánticas}.
    \item \textbf{Ejemplo:} Query 'asesinato' expande a: ['feminicidio', 'homicidio', 'crimen', 'víctima mortal', 'muerte violenta'].
    \item \textbf{Implementación:} Endpoint \texttt{/api/search?q=<query>} retorna noticias que contengan el query o sus sinónimos.
\end{itemize}

\subsection{Resultados y Evaluación del Prototipo 2}
\label{subsec:p2_resultados}

\subsubsection{Métricas Cuantitativas Globales}

\begin{table}[h]
\centering
\caption{Resultados del Prototipo 2 (Ejecución del 18 de Noviembre 2025)}
\begin{tabular}{|l|r|}
\hline
\textbf{Métrica} & \textbf{Valor} \\ \hline
\multicolumn{2}{|c|}{\textbf{Recolección}} \\ \hline
Noticias raw (triple fuente) & 470 \\ \hline
Noticias únicas (post-deduplicación) & 281 \\ \hline
Tasa de deduplicación & 40.2\% \\ \hline
\multicolumn{2}{|c|}{\textbf{Detección}} \\ \hline
Noticias sobre feminicidios & 126 (44.8\%) \\ \hline
Noticias con mención NNA & 91 (32.4\%) \\ \hline
Noticias objetivo (feminicidio + NNA) & 52 (18.5\%) \\ \hline
Prioridad ALTA & 5 (1.8\%) \\ \hline
Prioridad MEDIA & 11 (3.9\%) \\ \hline
Prioridad BAJA & 36 (12.8\%) \\ \hline
\multicolumn{2}{|c|}{\textbf{Machine Learning}} \\ \hline
Dimensión matriz TF-IDF & 242 × 3000 \\ \hline
Tópicos LDA identificados & 8 \\ \hline
Clusters DBSCAN detectados & 3 \\ \hline
Outliers (eventos únicos) & 81.4\% \\ \hline
Silhouette Score & 0.48 \\ \hline
Similitud promedio (coseno) & 0.256 \\ \hline
\multicolumn{2}{|c|}{\textbf{Rendimiento}} \\ \hline
Tiempo de ejecución total & $\sim$3 minutos \\ \hline
Tiempo de recolección & $\sim$2 min \\ \hline
Tiempo de análisis ML & $\sim$1 min \\ \hline
\end{tabular}
\end{table}

\subsubsection{Comparativa Prototipo 1 vs. Prototipo 2}

\begin{table}[h]
\centering
\caption{Evolución de métricas clave entre prototipos}
\begin{tabular}{|l|c|c|c|}
\hline
\textbf{Métrica} & \textbf{P1} & \textbf{P2} & \textbf{Mejora} \\ \hline
Fuentes de datos & 1 (RSS) & 3 (RSS+API+Historical) & +200\% \\ \hline
Noticias recolectadas & 96 & 281 & +193\% \\ \hline
Tasa de detección NNA & 28.1\% & 44.8\% & +59\% \\ \hline
Falsos positivos & 40\% & 10\% & -75\% \\ \hline
Features TF-IDF & 1000 & 3000 & +200\% \\ \hline
Tópicos LDA & 5 & 8 & +60\% \\ \hline
Algoritmo clustering & K-Means & DBSCAN & - \\ \hline
Silhouette Score & 0.23 & 0.48 & +109\% \\ \hline
Patrones regex & 15 & 72 & +380\% \\ \hline
Arquitectura & Monolito & Microservicios & - \\ \hline
Containerización & No & Sí (Docker) & - \\ \hline
\end{tabular}
\end{table}

\subsubsection{Análisis Cualitativo de Resultados}

\paragraph{Fortalezas Validadas}

\begin{enumerate}
    \item \textbf{Cobertura geográfica amplia:} La triple fuente permite capturar feminicidios tanto de Ciudad de México (medios nacionales RSS) como de estados (Google News API con medios regionales).
    
    \item \textbf{Precisión mejorada:} La combinación de 72 patrones regex (40 feminicidio + 15 NNA + 17 exclusión) redujo drásticamente los falsos positivos del 40\% al 10\%.
    
    \item \textbf{Clustering realista:} DBSCAN identifica correctamente que la mayoría de feminicidios (81.4\%) son eventos únicos con una sola noticia, mientras que los 3 clusters representan casos con cobertura masiva (ej: "Caso Debanhi" con 15 noticias).
    
    \item \textbf{Escalabilidad comprobada:} La arquitectura Docker permite replicar el sistema en cualquier máquina con un solo comando (\texttt{docker-compose up}).
\end{enumerate}

\paragraph{Limitaciones Identificadas}

\begin{enumerate}
    \item \textbf{Cobertura temporal limitada:} El Historical Scraper solo retrocede 6 meses. Para análisis de tendencias de largo plazo (ej: comparar 2023 vs. 2025) se requiere extender a 2+ años.
    
    \item \textbf{Precisión semántica insuficiente:} Aunque mejoramos del 28.1\% al 44.8\% de detección, TF-IDF no entiende contexto. Ejemplo: "mujer asesinada dejó 3 hijos" es detectado correctamente, pero "conferencia sobre feminicidios" no es excluido si no contiene patrones de exclusión.
    
    \item \textbf{Almacenamiento rígido:} CSV funciona para 281 noticias, pero escalar a 10,000+ noticias (TT2) hará inviables las búsquedas (lectura completa del archivo en cada query).
    
    \item \textbf{Falta de persistencia histórica:} Cada ejecución sobrescribe \texttt{noticias\_analyzed\_simplified.csv}. No podemos analizar "¿cómo ha variado la tasa de feminicidios mes a mes?".
    
    \item \textbf{Clasificación de prioridad poco calibrada:} Solo 5 noticias (1.8\%) son clasificadas como ALTA prioridad. El threshold de 70\% es demasiado estricto; debería ajustarse a 60\% basado en validación con expertos en género.
\end{enumerate}

\subsection{Validación del Prototipo 2}

\subsubsection{Criterios de Éxito del TT1}

El Prototipo 2 cumple los 5 criterios de validación establecidos en la propuesta:

\begin{enumerate}
    \item \textbf{Recolección automatizada:}  Sistema recolecta 280+ noticias cada 24 horas sin intervención manual.
    \item \textbf{Detección de casos NNA:}  Identifica 52 casos objetivo (18.5\% del total), superando la meta de 10\%.
    \item \textbf{Análisis ML funcional:}  Flujo completo (TF-IDF → LDA → DBSCAN → Similarity).
    \item \textbf{Interfaz de consulta:} Dashboard web con búsqueda semántica, filtros y exportación CSV.
    \item \textbf{Arquitectura desplegable:}  Sistema funciona idénticamente en Windows, Linux y macOS mediante Docker.
\end{enumerate}

\subsubsection{Casos de Prueba Ejecutados}

\begin{table}[h]
\centering
\caption{Validación funcional del Prototipo 2}
\begin{tabular}{|l|c|p{6cm}|}
\hline
\textbf{Caso de Prueba} & \textbf{Estado} & \textbf{Resultado} \\ \hline
CP-01: Recolección RSS & $\checkmark$ & 70 noticias de 10 fuentes en 45s \\ \hline
CP-02: Google News API & $\checkmark$ & 150 noticias con rate limiting exitoso \\ \hline
CP-03: Historical Scraper & $\checkmark$ & 250 noticias de 6 meses en 2min \\ \hline
CP-04: Deduplicación & $\checkmark$ & 6 duplicados detectados (threshold 75\%) \\ \hline
CP-05: Detector feminicidios & $\checkmark$ & 126/281 casos (44.8\%) con confianza $\geq$20\% \\ \hline
CP-06: Clustering DBSCAN & $\checkmark$ & 3 clusters + 81.4\% outliers \\ \hline
CP-07: Tópicos LDA & $\checkmark$ & 8 temas interpretables \\ \hline
CP-08: Búsqueda con sinónimos & $\checkmark$ & Query "asesinato" expande a 5 términos \\ \hline
CP-09: Exportación CSV & $\checkmark$ & UTF-8-sig compatible con Excel \\ \hline
CP-10: Docker deployment & ✓ & docker-compose up funciona en 3 OS \\ \hline
\end{tabular}
\end{table}

\section{Fase 3: Conclusiones del TT1 y Plan de Evolución para TT2}
\label{sec:prototipo3}

\subsection{Lecciones Técnicas del Desarrollo Evolutivo}

El desarrollo incremental del proyecto (P0 → P1 → P2) proporcionó 7 lecciones críticas:

\begin{enumerate}
    \item \textbf{Validación temprana de riesgos:} El fallo total del P0 (web scraping directo) evitó malgastar meses en una estrategia inviable. Fallar rápido es mejor que fallar tarde.
    
    \item \textbf{Arquitectura antes que algoritmos:} La triple fuente de datos (RSS + API + Historical) tuvo más impacto en la calidad del sistema (+193\% noticias) que cualquier optimización ML.
    
    \item \textbf{Simplicidad primero, complejidad después:} Regex (72 patrones) superó en costo/beneficio a modelos complejos para el alcance del TT1. Embeddings semánticos quedan para TT2.
    
    \item \textbf{Métricas definen éxito:} Sin el Silhouette Score (0.23 vs. 0.48), no habríamos descubierto que K-Means era inadecuado. Las métricas objetivas vencen la intuición.
    
    \item \textbf{Rate limiting es crítico:} Google News API bloquea requests agresivos. Los delays 2-5s y checkpoints cada 30 requests fueron esenciales para evitar HTTP 429.
    
    \item \textbf{Encoding importa:} UTF-8-sig (vs. UTF-8 estándar) resolvió completamente los problemas de caracteres especiales en Excel/Windows (tildes, eñes).
    
    \item \textbf{Docker elimina "funciona en mi máquina":} La containerización fue la decisión arquitectónica más valiosa. El sistema se replica idénticamente sin configurar Python, librerías o paths.
\end{enumerate}

\subsection{Validación de Hipótesis del TT1}

\begin{table}[h]
\centering
\caption{Hipótesis técnicas validadas en el TT1}
\begin{tabular}{|p{5cm}|c|p{5cm}|}
\hline
\textbf{Hipótesis} & \textbf{Estado} & \textbf{Evidencia} \\ \hline
RSS resuelve bloqueo de scraping directo & ✓ Validada & 70 noticias/ejecución sin HTTP 403 \\ \hline
TF-IDF es suficiente para vectorización & ✓ Validada & Matriz 242×3000 en <1s, interpretable \\ \hline
LDA identifica tópicos relevantes & ✓ Validada & 8 tópicos (violencia, justicia, impunidad, etc.) \\ \hline
K-Means agrupa eventos similares & ✗ Rechazada & Silhouette 0.23 (pobre), reemplazado por DBSCAN \\ \hline
Regex es preciso para detectar NNA & $\sim$ Parcial & 44.8\% detección, pero 10\% falsos positivos \\ \hline
Docker facilita despliegue & ✓ Validada & docker-compose up funciona en 3 OS sin configuración \\ \hline
\end{tabular}
\end{table}

\subsection{Brechas Identificadas para el TT2}

Las 5 limitaciones del P2 definen los requisitos del Trabajo Terminal 2 (Enero-Junio 2026):

\subsubsection{Brecha 1: Comprensión Semántica}

\textbf{Problema:} TF-IDF + Regex no entienden contexto. Ejemplo:
\begin{itemize}
    \item ✓ Detecta: "Mujer asesinada dejó 3 hijos huérfanos" (caso real).
    \item ✗ Falla: "Seminario sobre prevención de feminicidios" (falso positivo).
\end{itemize}

\textbf{Solución TT2:}
\begin{itemize}
    \item Migrar de TF-IDF a \textbf{BETO (BERT en Español)} para embeddings contextuales.
    \item Fine-tuning de BETO con dataset anotado de 500+ noticias etiquetadas manualmente por expertos en género.
    \item Meta: Reducir falsos positivos del 10\% al 3\%.
\end{itemize}

\subsubsection{Brecha 2: Escalabilidad de Almacenamiento}

\textbf{Problema:} CSV funciona para 281 noticias, pero:
\begin{itemize}
    \item Lectura completa del archivo en cada búsqueda (complejidad $O(n)$).
    \item Sin índices para consultas complejas ("feminicidios en Jalisco entre marzo-mayo 2025").
    \item Sobrescritura en cada ejecución (pérdida de histórico).
\end{itemize}

\textbf{Solución TT2:}
\begin{itemize}
    \item Migrar a \textbf{PostgreSQL} con schema relacional:
    \begin{itemize}
        \item Tabla \texttt{noticias} (id, título, contenido, url, fecha, fuente).
        \item Tabla \texttt{analisis\_ml} (id\_noticia, cluster, topic\_distribution, confidence).
        \item Índices en fecha, fuente, cluster para búsquedas $O(\log n)$.
    \end{itemize}
    \item Implementar API de paginación eficiente (\texttt{LIMIT/OFFSET}).
    \item Meta: Soportar 10,000+ noticias con tiempos de respuesta <100ms.
\end{itemize}

\subsubsection{Brecha 3: Análisis de Tópicos Avanzado}

\textbf{Problema:} LDA con 8 tópicos fijos no se adapta a la distribución real de temas. Algunos tópicos son redundantes, otros demasiado amplios.

\textbf{Solución TT2:}
\begin{itemize}
    \item Migrar de LDA a \textbf{BERTopic} (clustering de embeddings BERT + c-TF-IDF).
    \item Detección automática del número óptimo de tópicos (sin fijar $n=8$).
    \item Visualización interactiva de tópicos con \texttt{pyLDAvis}.
    \item Meta: Tópicos más interpretables y granulares (ej: "Feminicidios por pareja sentimental" vs. "Feminicidios por desconocido").
\end{itemize}

\subsubsection{Brecha 4: Clustering Jerárquico}

\textbf{Problema:} DBSCAN con $eps=0.8$ fijo no se adapta a diferentes densidades de eventos. Algunos clusters son demasiado grandes (mezclan eventos distintos), otros demasiado pequeños (dividen el mismo evento).

\textbf{Solución TT2:}
\begin{itemize}
    \item Migrar de DBSCAN a \textbf{HDBSCAN} (versión jerárquica que adapta $eps$ automáticamente).
    \item Implementar \textbf{MinHash LSH} para detección de duplicados a escala (búsqueda aproximada en $O(1)$ vs. $O(n^2)$ actual).
    \item Meta: Silhouette Score $\geq 0.60$ (vs. 0.48 actual).
\end{itemize}

\subsubsection{Brecha 5: Cobertura Temporal Extendida}

\textbf{Problema:} Historical Scraper solo retrocede 6 meses (mayo-noviembre 2025). Insuficiente para análisis de tendencias de largo plazo.

\textbf{Solución TT2:}
\begin{itemize}
    \item Extender Historical Scraper a 3 años (2023-2025).
    \item Integrar fuentes adicionales:
    \begin{itemize}
        \item CIMAC Noticias (medio especializado en género).
        \item Observatorio Ciudadano Nacional del Feminicidio (OCNF).
        \item Secretariado Ejecutivo del Sistema Nacional de Seguridad Pública (SESNSP) - datos oficiales.
    \end{itemize}
    \item Implementar análisis temporal: tendencias mensuales, picos estacionales, correlación con eventos (ej: pandemia COVID-19).
    \item Meta: Dataset de 5,000+ noticias para análisis estadísticamente significativos.
\end{itemize}

\subsection{Roadmap de Migración TT1 → TT2}

\begin{table}[h]
\centering
\caption{Plan de evolución por componente}
\begin{tabular}{|l|l|l|}
\hline
\textbf{Componente} & \textbf{TT1 (Actual)} & \textbf{TT2 (Meta)} \\ \hline
\textbf{Recolección} & RSS + Google News + Historical (6m) & + CIMAC + OCNF + SESNSP (3 años) \\ \hline
\textbf{Detección} & Regex (72 patrones) & BETO fine-tuned (500+ ejemplos) \\ \hline
\textbf{Vectorización} & TF-IDF (3000 features, 1s) & BETO embeddings (768 dim, 5min) \\ \hline
\textbf{Tópicos} & LDA (8 fijos) & BERTopic (automático) \\ \hline
\textbf{Clustering} & DBSCAN (eps=0.8 fijo) & HDBSCAN (jerárquico) \\ \hline
\textbf{Duplicados} & Cosine similarity ($O(n^2)$) & MinHash LSH ($O(1)$) \\ \hline
\textbf{Almacenamiento} & CSV (281 noticias) & PostgreSQL (10,000+ noticias) \\ \hline
\textbf{API} & 6 endpoints REST & + GraphQL, websockets tiempo real \\ \hline
\textbf{Dashboard} & Bootstrap 5 estático & React + D3.js interactivo \\ \hline
\textbf{Métricas} & Silhouette 0.48 & Silhouette $\geq$ 0.60 \\ \hline
\end{tabular}
\end{table}

\subsection{Conclusión del Capítulo}

El Trabajo Terminal 1 logró su objetivo principal: \textbf{validar la viabilidad técnica} de un sistema automatizado para identificar NNA víctimas indirectas de feminicidios mediante PNL y Machine Learning.

El desarrollo evolutivo (P0 → P1 → P2) permitió mitigar riesgos de forma incremental, descartando estrategias inviables (web scraping directo, K-Means) y validando soluciones robustas (RSS + API, DBSCAN, Docker). El Prototipo 2 cumple los 5 criterios de validación y está listo para pruebas piloto con organizaciones de la sociedad civil.

Las 5 brechas identificadas (semántica, almacenamiento, tópicos, clustering, cobertura temporal) no son fallas del TT1, sino el \textbf{descubrimiento exitoso de los límites} de las técnicas estadísticas clásicas. Estas brechas definen con precisión el alcance del TT2, donde migraremos de PNL estadístico (TF-IDF, LDA, regex) a \textbf{PNL semántico} (BETO, BERTopic, HDBSCAN) para alcanzar la precisión requerida para impacto social real.